\chapter{Multiple-Qubit Systems and Stream Cryptography (Suggested Problems 6)}

\noindent\textit{Sources: Eleanor Rieffel and Wolfgang Polak, \textit{Quantum Computing: A Gentle Introduction}, Chapter~3; Justin Wyss-Gallifent, \textit{Cryptography} (Chapter~13).}

\section{Why tensor products? (Rieffel \S3.1)}

Classically, combining two systems whose individual state spaces are $d$-dimensional often uses a \textbf{direct sum}, giving dimension $2d$ for two copies. Quantum mechanically, the joint state space of two \emph{distinct} quantum systems is modeled by the \textbf{tensor product} $V\otimes W$, with
\[
\dim(V\otimes W)=\dim(V)\cdot\dim(W).
\]
For $n$ qubits, each modeled by $\C^2$, the joint space is $(\C^2)^{\otimes n}$ and has dimension $2^n$---exponential in $n$. This ``bigger than classical'' state space is essential for entanglement and for quantum protocols.

\section{Direct sum vs.\ tensor product (\S3.1.1--3.1.2)}

\textbf{Direct sum $V\oplus W$:} If $\{\ket{\alpha_1},\ldots,\ket{\alpha_n}\}$ bases $V$ and $\{\ket{\beta_1},\ldots,\ket{\beta_m}\}$ bases $W$, then $V\oplus W$ has basis the \emph{disjoint union} (orthogonal direct sum) of these bases; $\dim(V\oplus W)=n+m$. Vectors are formal sums $\ket{v}\oplus\ket{w}$ with inner product $\braket{v_2\oplus w_2}{v_1\oplus w_1}=\braket{v_2}{v_1}+\braket{w_2}{w_1}$.

\textbf{Tensor product $V\otimes W$:} Bilinear construction generated by simple tensors $\ket{v}\otimes\ket{w}$, written $\ket{v}\ket{w}$. Bases $\{\ket{\alpha_i}\}$, $\{\ket{\beta_j}\}$ yield basis $\{\ket{\alpha_i}\otimes\ket{\beta_j}\}$ of size $nm$. Every element of $V\otimes W$ is a \emph{finite} sum of simple tensors, but most elements are \emph{not} a single simple tensor.

\textbf{Tensor product axioms (algebraic):} Bilinearity:
\[
(\ket{v_1}+\ket{v_2})\otimes\ket{w}=\ket{v_1}\otimes\ket{w}+\ket{v_2}\otimes\ket{w},\quad
\ket{v}\otimes(\ket{w_1}+\ket{w_2})=\cdots,\quad
(a\ket{v})\otimes\ket{w}=\ket{v}\otimes(a\ket{w})=a(\ket{v}\otimes\ket{w}).
\]

\textbf{Inner product on $V\otimes W$:} For inner product spaces,
\[
\braket{v_2\otimes w_2}{v_1\otimes w_1}=\braket{v_2}{v_1}\,\braket{w_2}{w_1},
\]
extended by linearity to sums (conjugate-linear in the first slot in the physics convention used by Rieffel). If $\{\ket{\alpha_i}\}$ and $\{\ket{\beta_j}\}$ are orthonormal, so is $\{\ket{\alpha_i}\otimes\ket{\beta_j}\}$.

\section{Two qubits: computational basis}

For single-qubit spaces each with standard basis $\{\ket{0},\ket{1}\}$, the two-qubit space has orthonormal \textbf{computational basis}
\[
\ket{00},\quad \ket{01},\quad \ket{10},\quad \ket{11},
\]
where $\ket{ab}=\ket{a}\otimes\ket{b}$. A \textbf{product state} has the form $\ket{\psi}=\ket{\phi}\otimes\ket{\chi}$. A state that cannot be written as a single product is \textbf{entangled}.

\subsection*{Worked example (Rieffel-style)}
If $\ket{v}=\tfrac{1}{\sqrt{5}}\begin{pmatrix}1\\-2\end{pmatrix}$ and $\ket{w}=\tfrac{1}{\sqrt{10}}\begin{pmatrix}-1\\3\end{pmatrix}$ in the standard basis of $\C^2$, then in \textbf{dictionary order} $(\ket{00},\ket{01},\ket{10},\ket{11})$,
\[
\ket{v}\otimes\ket{w}
=\tfrac{1}{5\sqrt{2}}\begin{pmatrix}-1\\3\\2\\-6\end{pmatrix},
\]
a unit vector in $\C^2\otimes\C^2\cong\C^4$. (This matches Rieffel Example~3.1.2; always fix an ordering of basis kets before identifying with column vectors.)

\subsection*{Counting dimensions}
If $\dim V=d$ and $\dim W=d'$, then $\dim(V\otimes W)=dd'$. For $n$ qubits, $\dim = 2^n$. The direct sum of $n$ copies of $\C^2$ would have dimension $2n$ only---wrong model for independent quantum subsystems.

\section{Entanglement: definitions and examples (\S3.1.2)}

\textbf{Separable vs.\ entangled.} A pure state in $V\otimes W$ is \textbf{separable} (or \textbf{product}) if it equals $\ket{v}\otimes\ket{w}$; otherwise it is \textbf{entangled}.

\textbf{Bell states (two qubits).} The \textbf{Bell basis} for $\C^2\otimes\C^2$ includes
\[
\ket{\Phi^+}=\tfrac{1}{\sqrt{2}}(\ket{00}+\ket{11}),\quad
\ket{\Phi^-}=\tfrac{1}{\sqrt{2}}(\ket{00}-\ket{11}),\quad
\ket{\Psi^\pm}=\tfrac{1}{\sqrt{2}}(\ket{01}\pm\ket{10}).
\]
These are maximally entangled for the bipartition into the two qubits.

\textbf{GHZ and W families ($n$ qubits).} For $n\ge 2$,
\[
\ket{\mathrm{GHZ}_n}=\tfrac{1}{\sqrt{2}}(\ket{0}^{\otimes n}+\ket{1}^{\otimes n}),\qquad
\ket{W_n}=\tfrac{1}{\sqrt{n}}(\ket{0\ldots 01}+\ket{0\ldots 10}+\cdots+\ket{10\ldots 0})
\]
(with exactly one $\ket{1}$ in each term). Both are entangled for $n\ge 2$ with respect to the decomposition into $n$ tensor factors (course problems: prove first for $n=2$ when the general statement is assigned for all $n>1$).

\section{Schmidt decomposition (two qubits, tool)}

Any $\ket{\psi}\in\C^2\otimes\C^2$ admits \textbf{Schmidt form} $\ket{\psi}=\sum_{k=1}^r \sqrt{\lambda_k}\ket{u_k}\otimes\ket{v_k}$ with $\{\ket{u_k}\}$, $\{\ket{v_k}\}$ orthonormal and $\lambda_k>0$. The \textbf{Schmidt rank} $r$ equals $1$ iff $\ket{\psi}$ is product; hence $r\ge 2$ implies entanglement.

\section{Suggested Problems 6 --- Rieffel alignment}

\noindent\textbf{Course handout:} For exercises asking to show a property for all $n>1$, it suffices to treat \textbf{$n=2$} (two-qubit systems) in this course.

\noindent\textbf{Quiz / Exam (same list):} Rieffel--Polak Exercises \textbf{3.1--3.6}.

\section{Modular arithmetic mod $2$ (Justin \S13.3)}

Bits live in $\mathbb{F}_2=\{0,1\}$ with addition and multiplication mod $2$. Equivalently, addition is \textbf{XOR} on bits: $0\oplus 0=0$, $0\oplus 1=1$, $1\oplus 1=0$. Matrices and vectors over $\mathbb{F}_2$ are used throughout; Gaussian elimination works mod $2$.

\section{Stream cipher: XOR with a shared key (\S13.4)}

Alice and Bob share a binary \textbf{key} $k_1k_2\ldots k_L$ (periodically repeated to match message length). \textbf{Plaintext} bits $p_i$ are encrypted as
\[
c_i \equiv p_i + k_i \pmod 2,
\]
\textbf{Ciphertext} $\{c_i\}$. Decryption is identical: $p_i\equiv c_i+k_i\pmod 2$ because $x+x\equiv 0\pmod 2$.

\section{Linearly recursive keys (\S13.5)}

\textbf{Definition (Justin 13.5.0.1).} A \textbf{linearly recursively defined key} of length $i$ is specified by
\begin{itemize}[nosep]
  \item \textbf{Initial bits} $\bar{s}=[x_1;\ldots;x_i]\in\{0,1\}^i$;
  \item \textbf{Coefficient vector} $\bar{c}=[1;c_2;\ldots;c_i]\in\{0,1\}^i$ (first entry $1$).
\end{itemize}
For $n\ge i+1$,
\[
x_n \equiv x_{n-i} + c_2 x_{n-i+1} + \cdots + c_i x_{n-1} \pmod 2,
\]
i.e.\ each new bit is a fixed $\mathbb{F}_2$-linear combination of the preceding $i$ bits. The pair is written $K=\{\bar{s},\bar{c}\}$.

\textbf{Periodicity.} Such a linear recurrence over $\mathbb{F}_2$ produces a \textbf{periodic} infinite key; the (shortest) period can be found by scanning for repeats (Justin discusses reversibility and determinants in \S13.6).

\section{Breaking the key (\S13.6--13.7)}

\textbf{Brute force:} Given a long \textbf{fragment} of the key, try smallest $i=2,3,\ldots$ and test all $\bar{c}\in\{0,1\}^i$ with $c_1=1$ until the recurrence is consistent with the fragment; then deduce period.

\textbf{Refined brute force (\S13.6.3):} Uses determinants of structured matrices built from the fragment to guess the shortest recurrence length before enumerating all coefficients.

\textbf{Linear algebra mod $2$ (\S13.7):} Solving for $\bar{c}$ reduces to linear systems $A\bar{c}\equiv \bar{b}\pmod 2$.

\section{Suggested Problems 6 --- Justin alignment}

\noindent\textbf{Quiz-level:} \S13.1--13.5.

\noindent\textbf{Exam-level (Justin only):} Complement of quiz within Chapter~13, i.e.\ \S13.6 onward (including exercises 13.6, 13.7, \ldots).

\noindent\textbf{Code:} \texttt{ex11\_sp6.m} (Octave: encryption, recursive keys, brute-force recursion checks).
